\documentclass[11pt]{article}
\usepackage[margin=2.6cm]{geometry}
\usepackage{amsmath,amssymb,amsthm}
\usepackage{graphicx}
\usepackage{booktabs}
\usepackage{hyperref}
\usepackage{caption}

\newtheorem{theorem}{Theorem}
\newtheorem{corollary}{Corollary}
\theoremstyle{definition}
\newtheorem{remark}{Remark}

\title{\bf The Operator That Almost Exists:\\ Eight New Campaigns and the Measured Boundary\\ of the Hilbert--P\'olya Program (v2.0)}
\author{Ishak Khamzatovich Isaev\\
\small Independent researcher, Nalchik, Kabardino-Balkarian Republic\\
\small \texttt{Hilbert Polya Bridge v2.0.0}}
\date{October 12, 2026}

\begin{document}
\maketitle

\begin{abstract}
Version 1.0 of the Hilbert--P\'olya Bridge measured the boundary between
the graph ($\xi$ on the critical line) and any lattice operator through
three projections of one gap (rank, prime staircase, heat trace) with
positive controls. Version 2.0 executes the v1 roadmap and adds eight
campaigns (C9--C16) with eight theorems (T11--T18). The headline results:
\textbf{(T11)} the heat trace with the Guinand--Weil kernel forces the
spectrum to be the zero multiset, so a prime-weight-in-trace primitive
exists and can be made arbitrarily nonlocal ($r_{H7}=1.000$,
$r_{EF}=0.999$) --- the measured H2 gap is a wall of \emph{locality};
\textbf{(T13)} the first arithmetic-to-spectrum bridge inside the lattice:
the $\alpha=1/3$ zero modes satisfy the counting law
$N_0 = 2L - 6\,[L \equiv 0 \pmod{12}]$ with Dirac index $-2$ on the
$L \equiv 0 \pmod{12}$ subfamily, verified through $L=48$;
\textbf{(T14)} the verified-RH band is extended to $T=5000$
($N = 4520 = N_{\mathrm{strip}}$, $S_{\max}=0.91$);
\textbf{(T17)} the trace formula transfers to Dirichlet $L$-functions with
the twisted weight $\chi(n)\Lambda(n)/\sqrt{n}$ verified to
$|\Delta K| \le 9.6\cdot 10^{-6}$, and the gates stay in the null zone
against $L$-zeros --- the H2 gap is \emph{universal} across the class.
Blind gate families (T15), a 121-configuration anomaly taxonomy (T16), and
a quantum-information layer (T18) confirm the demarcation. The package is
fully self-contained: the zeta zeros are computed and committed internally
(no external data files), verification is 18/18 IVP + 21/21 pytest, and an
open challenge with a public scoring protocol is attached.
\end{abstract}

\section{Introduction}
The Hilbert--P\'olya program asks whether the nontrivial zeros of $\zeta$
are the spectrum of a self-adjoint operator. The v1 package (C1--C8,
T1--T10, 87 gates) answered five questions by measurement, with positive
controls on machine precision. The roadmap named four directions able to
move the main gap: nonlocal trace constructions (the T6 lesson), stitching
the statistical and coding branches, a number-theoretic explanation of the
$\alpha = 1/3$ zero modes, and extending the verified band. v2.0 executes
all of them (C9, C10, C11, C12), adds blind campaigns and an anomaly hunt
(C13, C14), transfers the pipeline to Dirichlet $L$-functions (C15), and
adds a quantum-information layer (C16).

\section{C9/T11: the nonlocal trace primitive}
\begin{theorem}[Heat-trace uniqueness]
Let $H$ be self-adjoint with discrete spectrum $\{\lambda_n\} \subset
(0,\infty)$. If
\[
\operatorname{Tr} e^{-xH} \;=\; \frac{1}{\pi}\int_0^\infty Q(a)\sin(ax)\,
da \quad \text{for all } x>0,
\qquad Q(a) = \sum_{\gamma>0} \frac{2a}{a^2+\gamma^2},
\]
then $\operatorname{spec} H = \{\pm\gamma_n\}$ with multiplicity.
\end{theorem}
\begin{proof}
Both sides are Laplace--Stieltjes transforms of measures on $(0,\infty)$;
by uniqueness the measure is determined by the trace. Since
$\frac{1}{\pi}\int_0^\infty \frac{2a\,\sin(ax)}{a^2+\gamma^2}da =
e^{-x\gamma}$, the measure is the counting measure of the zero set.
\end{proof}
\begin{corollary}
The prime weight in the trace is the spectral measure itself: any operator
with such a trace is isospectral to $H_\xi$. Nonlocality is achievable
arbitrarily (dense orthogonal conjugations), locality is not.
\end{corollary}

Measured family: $H_\xi = \operatorname{diag}(\pm\gamma)$ (control); DCT
and QR twins $O\,\mathrm{diag}(\gamma)\,O^T$ (dense, isospectral to
$7\cdot10^{-13}$): $r_{H7} = 1.000$, $r_{EF}^{(\mathrm{fluct})} = 0.999$;
truncated Berry--Keating $xp$: $r_{H7}=0.047$ (null); the multiplicative
kernel $\Lambda(\mathrm{lcm}(m,n))/\sqrt{mn}$: degenerate spectrum (an
arithmetic signature), $r_{EF}=-0.96$. The T6 machine re-verifies on the
internal zero cache to $|\Delta K| \le 1.8\cdot10^{-6}$.

\section{C10/T12: hard co-location and its limit}
Any phase system with $\Phi_{pq} \equiv \Phi_p + \Phi_q \pmod{2\pi}$ is a
homomorphism from the free abelian group on the primes; its normal form is
$\Phi_p = 2\pi\{\beta \ln p\}$. The $\{0,\pi\}$-valued case
($\Phi_p = \pi\chi_4(p)$) is \emph{gauge-trivial} in the lattice overlay
channel: $e^{\pm i\pi} = -1$ flips every $y$-bond and $\psi \to
(-1)^y\psi$ removes it (spectra coincide to $8.4\cdot10^{-15}$). With the
site-linked overlay the hard co-location holds at machine zero
($4.4\cdot10^{-16}$) for every $\beta$, yet the spectrum does not see it:
$r_{H7} \le 0.092$ over an 11-point $\beta$-grid (Bonferroni $0.20$), GUE
class unchanged. Universality erases even hard codings.

\section{C11/T13: the arithmetic bridge at $\alpha = 1/3$}
For the K\"ahler--Dirac (staggered) operator in the $\gamma_5$-Hermitian
form $H = \varepsilon D + m$ with Landau flux $\alpha = 1/3$ on the torus
($\alpha L \in \mathbb{Z}$ forces $L \equiv 0 \pmod 3$), the full scan
$L = 6,\dots,48$ gives the counting law
\[
N_0(L) \;=\; 2L \;-\; 6\,\big[L \equiv 0 \pmod{12}\big]
\quad\text{for } L \equiv 0 \pmod 6, \qquad N_0 = 0 \text{ otherwise},
\]
with Dirac index $\mathrm{ind} = \sum_{\text{zero modes}}
\mathrm{sign}(\varepsilon) = -2\,[L \equiv 0 \pmod{12}]$. The decisive
checks: $L = 36 \to 66$, $L = 48 \to 90$ (both exact); open boundaries
give one edge zero mode with index $-1$ for odd $L \equiv 0 \pmod 3$; the
$\alpha = 1/2$ control gives none. The v1 observation «tori $L \equiv 0
\pmod 4$» described only the index $-2$ subfamily: $L = 18$ ($\equiv 2
\pmod 4$) also carries $36$ modes. This is the first
arithmetic-to-spectrum bridge of the program; proving the law through the
K\"ahler--Dirac index is the v3 goal.

\section{C12/T14: the band to $T = 5000$}
With $Z(t) = \mathtt{mp.siegelz}$ (bridge to $e^{i\theta}\zeta$: max
$|\Delta| = 0$ at six heights), adaptive sign-change counting (coarse grid
$0.15$, small-$|Z|$ refinement, 36-iteration bisection) and Turing-monitored
$S(T)$, the count gives $N_{\mathrm{online}}(5000) = 4520$ against
$\bar N(5000) = 4520.08$, $S_{\max} = 0.91 \le 2$; three full-strip
argument-principle windows match ($11{=}11$, $10{=}10$, $5{=}5$); the v1
anchors $(0,3,10,29,79,649)$ reproduce bit-wise. Gaps below 5000: min
$0.0908$, mean $1.103$.

\section{C13/T15 + C14/T16: blind families and the attractor taxonomy}
Five families absent from the v1 handle space (imaginary $x$-hoppings;
golden-ratio quasiperiodic vortices; vortex dimers; power-law hoppings
$r^{-2}$; row flux dithering) produce 20 gates: all in the null zone
($|r_{H7}| \le 0.040$, $\min p = 0.29$) with the pipeline control
$PC0' = 1.000$. A 121-configuration scan yields the honest taxonomy
$97$ GUE $/\,18$ integrable ($N_v{=}0$, Poisson) $/\,6$ intermediate
(weak-vortex $N_v \le 8$, $q{=}0.5$, $\langle r\rangle = 0.47$--$0.53$);
no rank outliers ($|r_{H7}| \le 0.153$). The GUE attractor is global for
vortex-carrying configurations, with a measurable basin structure.

\section{C15/T17: Dirichlet $L$-functions — the gap is universal}
For the quadratic characters $\chi_5$ (even) and $\chi_7$ (odd):
$Z_\chi(t) = e^{i\theta_\chi(t)}L(\tfrac12{+}it,\chi)$ is real to
$5\cdot10^{-35}$; bisection finds $68$ and $74$ zeros $\le 120$ against
smooth counts $67.8$ and $74.5$; the twisted staircase
$\psi_\chi(x) = \sum_{n\le x}\chi(n)\Lambda(n)$ matches the zero reference
at $0.902$/$0.949$.
\begin{theorem}[Twisted trace formula]
$K_\chi(x) = \sum_\gamma e^{-x\gamma_\chi} = \frac{1}{\pi}\int_0^\infty
Q_\chi(a)\sin(ax)\,da$, where $Q_\chi(a) = \frac{\Lambda'}{\Lambda}
(\tfrac12{+}a,\chi) - B_\chi - C_0^{(\chi)}$ carries the twisted weight
$\chi(n)\Lambda(n)/\sqrt{n}$, and $C_0^{(\chi)} =
\frac{\Lambda'}{\Lambda}(\tfrac12) - B_\chi$.
\end{theorem}
Verified to $|\Delta K| \le 9.6\cdot10^{-6}$ ($\chi_5$) and
$6.1\cdot10^{-6}$ ($\chi_7$). The same lattice gates against $L$-zeros:
$r_{H7} = 0.035$--$0.164$, the null zone. The H2 gap is a class property,
not a zeta defect.

\section{C16/T18: the quantum-information layer}
Fermionic entanglement of the Slater Fermi sea (correlation-matrix
protocol): the gate's half-system entropy $S_A \approx 12$--$16$ nats
against the Haar baseline $160$--$220$ ($z>100$) --- locality suppresses
the volume-law entanglement, the quantum echo of C9/T11. The mutual
information profile carries no zero signature ($-0.018$ vs $-0.121$
random). Spectral-form-factor slopes on unfolded 200-level windows are
noisy; the gate matches GUE within the spread.

\section{Honesty ledger (found-and-fixed, v2)}
(i) the shipped $B_{\mathrm{Hadamard}}$ constant used $\ln 2\pi$ instead
of $\ln 4\pi$ (committed v1 verdicts unaffected);
(ii) the v1 T9 observation refined to the full T13 law;
(iii) $\{0,\pi\}$-valued codings are gauge-trivial (T12 corollary);
(iv) the $r_{EF}$ metric is taken on detrended fluctuations (raw
correlation catches common trends); plus the T6-check fixes ($n$ vs
$\ln n$; the $a=0.5$ pole excluded from the spline grid).

\section{Reproducibility and the open challenge}
The package uses no external data files: the zeta zeros are computed by
mpmath internally and committed (800 ordinates to $\gamma = 1183.7$,
bit-reproducible). Verification: IVP01--IVP18 (18/18 PASS) and pytest
(21/21). An open challenge is attached: demonstrate $r_{H7} \ge 0.7$
under the v2 protocol with passing positive controls and an intact GUE
class; a blind scorer (\texttt{challenge/score\_submission.py}) rebuilds
the gate from the submitted spec and archives the full report.

\section{The v2 verdict}
The operator that almost exists exists exactly as far as the trace
allows: with the prime weight in it (T11), with arbitrary nonlocality,
with prime co-location in its phases (T12). It refuses to exist locally:
none of 227 measured configurations of five families and three classes
left the null zones (T15, T16). Arithmetic entered the lattice spectrum
where nobody looked: in the Dirac-index zero-mode count (T13). And the
gap measured on zeta is a property of the whole $L$-function class
(T17). The continuation is the v3 roadmap: proving the T13 law through
the index, other characters and modular forms, stitching through the
index, the band to $T = 50000$, and the quantum response to the index.

\begin{thebibliography}{9}
\bibitem{riemann1859} B. Riemann (1859), \emph{On the Number of Primes
Less Than a Given Magnitude}.
\bibitem{hadamard1893} J. Hadamard (1893), canonical product;
the constant $B = \xi'(0)/\xi(0)$.
\bibitem{weil1952} A. Weil (1952), explicit formula for test functions;
Guinand's zero-side formulation.
\bibitem{bgs1984} O. Bohigas, M.-J. Giannoni, C. Schmit (1984), GUE
hypothesis for quantum spectra.
\bibitem{berrykeating1999} M. Berry, J. Keating (1999), $H = xp$ and the
semiclassical window.
\bibitem{connes1999} A. Connes (1999), trace formula and the
semilocal factor.
\bibitem{hpbv1} I. K. Isaev (2026), \emph{Hilbert--P\'olya Bridge v1.0.0:
a self-contained numerical study of the graph--operator passage},
this repository.
\end{thebibliography}

\end{document}
